% !TEX program = xelatex
% ============================================================
%  الرد على المراجعين — Response to Reviewers Template
%  نموذج رسالة الرد على مراجعات المجلة العلمية (Point-by-Point Response)
%  Compile with: xelatex main.tex (run twice)
% ============================================================
%  Features:
%    - Letter format with editor name and manuscript ID
%    - Point-by-point response table (Reviewer comment | Author response | Line numbers)
%    - tabularx + booktabs for clean response tables
%    - Separate blocks per reviewer
%    - Fancy headers with manuscript title
% ============================================================

\documentclass[11pt,a4paper]{article}

% ---- Geometry ----
\usepackage[a4paper,margin=2.5cm,top=2.5cm,bottom=2.5cm,headheight=16pt]{geometry}

% ---- Core packages ----
\usepackage{graphicx}
\usepackage{array}
\usepackage{booktabs}
\usepackage{tabularx}
\usepackage{multirow}
\usepackage{enumitem}
\usepackage{float}
\usepackage{caption}
\usepackage{titlesec}
\usepackage{fancyhdr}
\usepackage{setspace}
\usepackage{xcolor}
\usepackage{amsmath}

% ---- RTL / Arabic language support (load before hyperref) ----
\usepackage{bidi}
\usepackage[no-math]{fontspec}
\usepackage[quiet]{polyglossia}

% ---- Fonts ----
\setmainfont[Scale=1]{Amiri-Regular.ttf}
\newfontfamily\arabicfont[Script=Arabic,Scale=0.95]{Amiri-Regular.ttf}
\newfontfamily\englishfont[Scale=0.95]{TeX Gyre Termes}
\setmonofont{Amiri-Regular.ttf}
\newfontfamily\arabicfontsf{Amiri-Regular.ttf}

% ---- Languages ----
\setdefaultlanguage[calendar=gregorian,hijricorrection=1]{arabic}
\setotherlanguage[variant=british]{english}

% ---- Hyperref ----
\usepackage[breaklinks,hidelinks]{hyperref}
\hypersetup{
  pdftitle={الرد على المراجعين},
  pdfauthor={اسم المؤلف}
}

% ---- Captions in Arabic ----
\addto\captionsarabic{%
  \renewcommand{\tablename}{الجدول}%
}

% ---- Section formatting ----
\titleformat{\section}{\Large\bfseries}{\thesection}{0.7em}{}
\titlespacing*{\section}{0pt}{1.4ex plus 0.4ex}{0.9ex plus 0.3ex}

% ---- Headers / footers ----
\pagestyle{fancy}
\fancyhf{}
\fancyhead[R]{\small\itshape الرد على المراجعين}
\fancyhead[L]{\small اسم المؤلف}
\fancyfoot[C]{\small\thepage}
\renewcommand{\headrulewidth}{0.3pt}
\renewcommand{\footrulewidth}{0pt}

% ---- List spacing ----
\setlist[itemize]{topsep=0.3em, itemsep=0.15em, parsep=0pt}
\setlist[enumerate]{topsep=0.3em, itemsep=0.15em, parsep=0pt}

% ---- Table column types ----
\newcolumntype{L}[1]{>{\raggedright\arraybackslash}p{#1}}
\newcolumntype{C}[1]{>{\centering\arraybackslash}p{#1}}
\newcolumntype{R}[1]{>{\raggedleft\arraybackslash}p{#1}}
\newcolumntype{Y}{>{\raggedright\arraybackslash}X}

% ---- Line spacing ----
\setstretch{1.4}

\begin{document}

% ============================================================
% LETTER HEADER
% ============================================================
\thispagestyle{fancy}

\begin{flushleft}
\textbf{\today}
\end{flushleft}

\vspace{1em}

\noindent
\textbf{إلى المحرر الموقر،}\\[0.5em]
% TODO: اكتب اسم المحرر هنا.
د. اسم المحرر\\
% TODO: اكتب اسم المجلة هنا.
اسم المجلة العلمية

\vspace{1em}

\noindent
\textbf{الموضوع:} الرد على مراجعات المخطوطة (Manuscript) رقم \\
% TODO: اكتب رقم المخطوطة هنا.
\textbf{\LR{MANUSCRIPT-ID-12345}}\\
% TODO: اكتب عنوان المخطوطة هنا.
\textbf{عنوان المخطوطة (Manuscript Title)}

\vspace{1em}

\noindent
المحرر الموقر،

\vspace{0.5em}

% TODO: اكتب مقدمة الرسالة هنا. اشكر المحرر والمراجعين على ملاحظاتهم.
نشكركم على مراجعة مخطوطتنا وإتاحة الفرصة لتحسينها. نقدم فيما يلي رداً تفصيلياً على ملاحظات المراجعين واحداً تلو الآخر (Point-by-Point). تظهر تعليقات المراجعين متبوعة بردنا مع الإشارة إلى أرقام الأسطر (Line numbers) المعدّلة في النسخة المنقحة.

\vspace{1em}

% ============================================================
% REVIEWER 1
% ============================================================
\section{الرد على ملاحظات المراجع الأول (Reviewer 1)}

% TODO: أضف جدول الرد على ملاحظات المراجع الأول.

\begin{table}[H]
\centering
\small
\begin{tabularx}{\textwidth}{L{5cm} Y L{2.5cm}}
\toprule
\textbf{تعليق المراجع} & \textbf{رد المؤلف} & \textbf{أرقام الأسطر} \\
\midrule
% TODO: انسخ تعليق المراجع الأول هنا.
\LR{Comment 1:} نص تعليق المراجع الأول هنا.
& % TODO: اكتب ردك هنا.
رد المؤلف على التعليق الأول مع شرح التعديل المُجرى.
& \LR{12--18} \\
\midrule
% TODO: انسخ تعليق المراجع الثاني هنا.
\LR{Comment 2:} نص تعليق المراجع الثاني هنا.
& % TODO: اكتب ردك هنا.
رد المؤلف على التعليق الثاني مع شرح التعديل المُجرى.
& \LR{45--52} \\
\bottomrule
\end{tabularx}
\end{table}

% ============================================================
% REVIEWER 2
% ============================================================
\section{الرد على ملاحظات المراجع الثاني (Reviewer 2)}

% TODO: أضف جدول الرد على ملاحظات المراجع الثاني.

\begin{table}[H]
\centering
\small
\begin{tabularx}{\textwidth}{L{5cm} Y L{2.5cm}}
\toprule
\textbf{تعليق المراجع} & \textbf{رد المؤلف} & \textbf{أرقام الأسطر} \\
\midrule
% TODO: انسخ تعليق المراجع الأول هنا.
\LR{Comment 1:} نص تعليق المراجع الأول هنا.
& % TODO: اكتب ردك هنا.
رد المؤلف على التعليق الأول مع شرح التعديل المُجرى.
& \LR{67--73} \\
\midrule
% TODO: انسخ تعليق المراجع الثاني هنا.
\LR{Comment 2:} نص تعليق المراجع الثاني هنا.
& % TODO: اكتب ردك هنا.
رد المؤلف على التعليق الثاني مع شرح التعديل المُجرى.
& \LR{98--105} \\
\bottomrule
\end{tabularx}
\end{table}

% ============================================================
% REVIEWER 3 (optional)
% ============================================================
\section{الرد على ملاحظات المراجع الثالث (Reviewer 3)}

% TODO: إذا كان هناك مراجع ثالث، أضف جدول الرد هنا. وإلا احذف هذا القسم.

\begin{table}[H]
\centering
\small
\begin{tabularx}{\textwidth}{L{5cm} Y L{2.5cm}}
\toprule
\textbf{تعليق المراجع} & \textbf{رد المؤلف} & \textbf{أرقام الأسطر} \\
\midrule
% TODO: انسخ تعليق المراجع الأول هنا.
\LR{Comment 1:} نص تعليق المراجع الأول هنا.
& % TODO: اكتب ردك هنا.
رد المؤلف على التعليق الأول مع شرح التعديل المُجرى.
& \LR{130--138} \\
\bottomrule
\end{tabularx}
\end{table}

% ============================================================
% CLOSING
% ============================================================
\section{ملاحظات ختامية}

% TODO: اكتب ملاحظات ختامية هنا.
نأمل أن تكون التعديلات المُجرىة قد عالجت جميع ملاحظات المراجعين. ونحن ممتنون لوقتهم وملاحظاتهم القيّمة التي ساهمت في تحسين جودة المخطوطة. نرحب بأي ملاحظات إضافية.

\vspace{1.5em}

\noindent
مع خالص التقدير،\\[1em]
% TODO: اكتب اسم المؤلف ومعلومات التواصل هنا.
\textbf{اسم المؤلف}\\
القسم، الجامعة\\
\texttt{email@example.com}

\end{document}
